
\subsubsection{5.1 Main Results}

Table 1 presents results on the TruthfulQA test set.

\textbf{Table 1: Results on TruthfulQA Test Set}

\begin{table}[htbp]
\centering
\resizebox{\columnwidth}{!}{%
\begin{tabular}{llll}
\toprule
Method & Spearman rho & p-value & AUROC \\
\midrule
SEP (baseline) & 0.0835 & 0.288 & 0.5214 \\
SEDP (proposed) & 0.0843 & 0.283 & 0.5219 \\
Delta (SEDP - SEP) & +0.0009 & --- & +0.0004 \\
Threshold (existence proof) & >= 0.30 & < 0.05 & --- \\
Random baseline & 0.00 & --- & 0.50 \\
\bottomrule
\end{tabular}
}
\end{table}

The existence proof threshold was not met. SEDP achieved rho = 0.0843, falling short of the 0.3 threshold by 72\%. The p-value of 0.283 indicates that the correlation is not statistically significant; the null hypothesis that the true correlation is zero cannot be rejected.

Both methods achieved AUROC approximately 0.52, marginally above the random baseline of 0.50. For practical hallucination detection, this performance provides no useful signal.

SEDP marginally outperformed SEP (+0.0009 rho, +0.0004 AUROC). While the effect direction is consistent with the hypothesis that similarity features provide additional information, the magnitude is negligible.

\subsubsection{5.2 Comparison with Published Results}

\textbf{Table 2: Comparison with Published Benchmarks}

\begin{table}[htbp]
\centering
\resizebox{\columnwidth}{!}{%
\begin{tabular}{ll}
\toprule
Source & AUROC on TruthfulQA \\
\midrule
Kossen et al. (2024) & ~0.85 \\
This work (SEP baseline) & 0.52 \\
Gap & 39\% \\
\bottomrule
\end{tabular}
}
\end{table}

The 39\% gap between the achieved AUROC and published results is substantial. The experiment followed published guidance: layer 25 falls within the recommended middle-to-late range, TBG is a standard token position, and logistic regression matches the published probe architecture.

Possible explanations for this gap include:

\begin{enumerate}
\item \textbf{Layer selection:} Layer 25 may not be optimal for Llama-3-8B-Instruct; systematic ablation across layers 20-31 was not performed.
\end{enumerate}

\begin{enumerate}
\item \textbf{Token position:} Alternative positions (SLT, pooling) may capture different information.
\end{enumerate}

\begin{enumerate}
\item \textbf{Implementation differences:} Subtle variations in response generation, SE computation, or preprocessing may contribute to the gap.
\end{enumerate}

\begin{enumerate}
\item \textbf{Model-specific factors:} Published results may use different model versions or configurations.
\end{enumerate}

\subsubsection{5.3 Effect of Similarity Augmentation}

Despite the overall failure, the effect direction was examined:

\begin{table}[htbp]
\centering
\resizebox{\columnwidth}{!}{%
\begin{tabular}{llll}
\toprule
Metric & SEP & SEDP & Delta \\
\midrule
Spearman rho & 0.0835 & 0.0843 & +0.0009 \\
AUROC & 0.5214 & 0.5219 & +0.0004 \\
\bottomrule
\end{tabular}
}
\end{table}

The positive delta is consistent with the hypothesis that similarity features contribute information. However, the magnitude is less than 1\% relative improvement, providing no practical benefit. When the base signal is absent, auxiliary features cannot compensate.

\subsubsection{5.4 Visualization}

Figure 1 presents the Spearman correlation for both methods with the existence proof threshold.

\begin{figure}[htbp]
\centering
\includegraphics[width=\columnwidth]{figures/gate_metrics.png}
\caption{Gate Metrics}
\label{fig:gate_metrics}
\end{figure}

\textit{Figure 1: Spearman correlation (rho) for SEP and SEDP. The horizontal dashed line indicates the existence proof threshold (rho = 0.3). Both methods achieved rho approximately 0.08, far below the threshold.}

Figure 2 shows the scatter plot of predicted versus true SE values.

\begin{figure}[htbp]
\centering
\includegraphics[width=\columnwidth]{figures/scatter.png}
\caption{Scatter Plot}
\label{fig:scatter}
\end{figure}

\textit{Figure 2: Predicted SE probability versus true semantic entropy for both methods. No discernible correlation is visible.}

Figure 3 presents ROC curves for both methods.

\begin{figure}[htbp]
\centering
\includegraphics[width=\columnwidth]{figures/roc_curves.png}
\caption{ROC Curves}
\label{fig:roc_curves}
\end{figure}

\textit{Figure 3: ROC curves for SEP and SEDP. Both curves lie close to the diagonal, indicating near-random classification performance (AUROC approximately 0.52).}
